% ============================================================
\section{Introduction}
\label{sec:intro}

Fault-tolerant quantum computation protects each logical qubit inside a patch of
a quantum error-correcting code, most commonly the surface code~\cite{fowler2012surface}.
Patches are then laid out on a two-dimensional grid~\cite{litinski2019game}. On such
hardware, a logical two-qubit gate is performed via \emph{lattice surgery}, which 
temporarily merges and then splits a strip of ancilla patches between the 
two operands~\cite{horsman2012lattice,litinski2019game}. 
Single-qubit Clifford gates carry no such cost, but the non-Clifford $T$ gate 
that completes a universal gate set does; 
it must be teleported by consuming a \emph{magic state}~\cite{bravyi2005magic} 
from a magic-state-factory patch~\cite{litinski2019game}, which must be connected to 
the operand, also through lattice surgery.

Compiling a logical circuit onto this fabric splits into two sub-problems. 
Placement (\emph{mapping}) assigns each logical qubit to a patch; \emph{routing}
then schedules the lattice-surgery operations, reserving disjoint
ancilla corridors at every time-step.
The initial placement is decisive: $T$-heavy qubits stranded far from a
magic-state factory, or frequently-interacting partners scattered across the grid,
often cause lattice surgeries to serialize and inflate routing depth.

We address the placement sub-problem with a \emph{potential-field} heuristic. 
Each candidate tile is scored by superposing two-dimensional Gaussians---attractive
toward magic-state factories and frequently-interacting partners, repulsive around
already-placed patches to keep ancilla corridors open. Every qubit is committed to
the highest-scoring tile that passes a \emph{safe-passage} test---a fast reachability
check guaranteeing that no placement traps a patch from reaching its partners or
magic-state factories---so the resulting mapping is always routable. 
The whole pass is deterministic and runs in
low-order polynomial time, in contrast to the heavier fault-tolerant compilation
strategies such as dependency-aware optimization, topology-centric
representations, or architecture-level co-optimization.

\smallskip
\noindent\textbf{Related work.}
Qubit mapping and routing are well studied in the NISQ regime~\cite{siraichi2018qubit}, 
where heuristics such as SABRE~\cite{li2019sabre} and the Qiskit transpiler~\cite{qiskit} 
insert $\SWAP{}$ gates to satisfy a fixed coupling graph; these methods
minimize $\SWAP{}$ count or depth, objectives that do not transfer to
our setting. Fault-tolerant compilation has since diversified into several
complementary lines. Its foundations are the surface-code and lattice-surgery
frameworks~\cite{fowler2012surface,horsman2012lattice} and Litinski's
\emph{game of surface codes}~\cite{litinski2019game}, whose tile-based space-time
model fixes the lattice-surgery cost model and tile budgets we optimize against;
optimizing lattice surgery exactly is itself $\NP$-hard~\cite{herr2017nphard}.
Routing-oriented methods recast long-range logical communication as graph
problems---edge-disjoint paths~\cite{beverland2022surface} and three-dimensional
path scheduling~\cite{hamada2026routing3d}---while end-to-end compilers and joint
optimizers attack mapping and routing together through scalable intermediate
representations~\cite{watkins2024compiler}, dependency-aware
search~\cite{wisq}, or resource-adaptive scheduling~\cite{zhu2024ecmas}. A
further line broadens the design space beyond conventional place-and-route: SAT
synthesis of volume-optimal subroutines such as $T$ 
factories~\cite{tan2024satscalpel}, multi-qubit lattice-surgery
scheduling~\cite{silva2024multiqubit}, topology-centric intermediate
representations based on ZX diagrams with search-based layout
synthesis~\cite{zhou2026topols}, and architecture-level co-optimization of compute
regions, routing space and magic-state
resources~\cite{kan2025sparo,bharadwaj2026qomet}. A recurring message of this
recent work is that no single compilation family is universally optimal: the
preferred strategy depends on workload properties such as circuit density,
logical-qubit count and $T$ fraction~\cite{leblond2025comparison}.

Our closest point of comparison is WISQ~\cite{wisq}, a recent end-to-end compiler
that treats placement and scheduling jointly with a dependency-aware solver; it
produces high-quality schedules but is far slower to compile. Relative to all of
these, our contribution is deliberately positioned in a different part of the
design space. Rather than solving a global mapping-and-routing problem online, we
compute a cheap, interpretable placement up front---encoding the dominant
geometric pressures of lattice-surgery execution---and hand it to a simple
downstream router. The method is thus complementary to global, topology-aware and
architecture-aware optimizers~\cite{wisq,kan2025sparo,zhou2026topols}: it is a
fast initial-placement strategy, not a replacement for end-to-end fault-tolerant
compilation flows.

\smallskip
\noindent\textbf{Contributions.}
We formulate fault-tolerant initial placement as the maximization of a Gaussian
potential field over the lattice, sourced by magic-state factories, mapped interaction
partners and already-placed qubits, and we pair it with a safe-passage filter
that certifies routability before any routing is attempted. We give two
weightings of the field---a discrete \emph{coarse} variant and a continuous
\emph{fine} variant that interpolates field strength from per-qubit $T$-count and
per-pair CNOT statistics---and implement everything in an open-source C++20
OpenQASM-to-lattice-surgery compiler~\cite{ftqc}. On 256 arithmetic,
algorithmic and state-preparation circuits the field placement reduces routing
steps by \todo{$X\%$} on average over a random baseline and matches
or improves on WISQ (by \todo{$M_{low}-M_{high}\times$}) for \todo{$k$ of 256} circuits, 
while compiling \todo{$Y\times$} faster (Section~\ref{sec:eval}).
